% GENERATED FILE. Edit canonical lessons, then run npm run generate:books.
% canonical-source-hash: fnv1a64:4ba2a3d3744aa223
% canonical-lessons: GE-C393-datei, GE-C393-speichern, GE-C393-ordner, GE-C393-anhang, GE-C393-anhaengen, GE-R393-review-386-393

\chapter{Dateien und Ordner}
\label{ch:ge-dateien-und-ordner}
\hlchaptermodality{393}

\begin{chapteropening}
\textbf{By the end of this chapter:} \emph{I can talk about files and folders, saving work, and sending attachments.}

Complete the last lesson of chapter 393: I can talk about files and folders, saving work, and sending attachments.
\end{chapteropening}

\section[die Datei]{die Datei — the (computer) file}

\label{lesson:GE-C393-datei}



\begin{quote}
Before the new one: say the German for the charger, then the German for the plug.
\end{quote}

\subsection*{You'll want to know: die Datei}
\textbf{die Datei} — \textquotedblleft{}the (computer) file\textquotedblright{}.

\begin{quote}
Die Datei ist zu groß für eine E-Mail.
\end{quote}

\begin{quote}
\emph{The file is too big for an email.}
\end{quote}

Plural: \emph{die Dateien}. \emph{Eine Datei öffnen} — to open a file.

\begin{cousinweb}
\emph{Datei} is usually explained as a 20th-century German blend of \emph{Daten} \textquotedblleft{}data\textquotedblright{}, from Latin \emph{datum} \textquotedblleft{}something given\textquotedblright{} (EtymDB), and \emph{Kartei} \textquotedblleft{}card index\textquotedblright{}. EtymDB has no entry for \emph{Datei} itself.
\end{cousinweb}

\begin{grammarlens}[title={What you've built}]
One more word for telling, explaining and describing at B1.
\end{grammarlens}

\subsection*{Your turn}
\begin{itemize}
\raggedright
  \item \emph{Say it:} \emph{die Datei}
  \item \emph{Say it:} \emph{die Datei}, once more
  \item \emph{Say it:} a sentence of your own with \emph{die Datei}
  \item \emph{Recall:} say the German for the charger, then the German for the plug, then say \emph{die Datei} again
\end{itemize}

\begin{tcolorbox}[breakable,colback=teal!4,colframe=teal!35!black,title={Before you move on}]
What does die Datei mean? (\textquotedblleft{}The (computer) file\textquotedblright{}.) Say it once more.
\end{tcolorbox}
\section[speichern]{speichern — to save (a file); to store}

\label{lesson:GE-C393-speichern}



\begin{quote}
Before the new one: say the German for the plug, then the German for the.
\end{quote}

\subsection*{You'll want to know: speichern}
\textbf{speichern} — \textquotedblleft{}to save (a file); to store\textquotedblright{}.

\begin{quote}
Ich habe vergessen, den Text zu speichern.
\end{quote}

\begin{quote}
\emph{I forgot to save the text.}
\end{quote}

\emph{Eine Datei speichern} — to save a file. \emph{Der Speicher} — the memory; also a storehouse or an attic.

\begin{cousinweb}
\emph{speichern} is made from \emph{Speicher} \textquotedblleft{}storehouse\textquotedblright{} (EtymDB), from Middle High German \emph{spīcher}, Old High German \emph{spīhhāri} and Late Latin \emph{spīcārium} \textquotedblleft{}a granary\textquotedblright{}, from \emph{spīca} \textquotedblleft{}an ear of grain\textquotedblright{}.
\end{cousinweb}

\begin{grammarlens}[title={What you've built}]
One more word for telling, explaining and describing at B1.
\end{grammarlens}

\subsection*{Your turn}
\begin{itemize}
\raggedright
  \item \emph{Say it:} \emph{speichern}
  \item \emph{Say it:} \emph{speichern}, once more
  \item \emph{Say it:} a sentence of your own with \emph{speichern}
  \item \emph{Recall:} say the German for the plug, then the German for the, then say \emph{speichern} again
\end{itemize}

\begin{tcolorbox}[breakable,colback=teal!4,colframe=teal!35!black,title={Before you move on}]
What does speichern mean? (\textquotedblleft{}To save (a file); to store\textquotedblright{}.) Say it once more.
\end{tcolorbox}
\section[der Ordner]{der Ordner — the folder; the ring binder}

\label{lesson:GE-C393-ordner}



\begin{quote}
Before the new one: say the German for the, then the German for to save.
\end{quote}

\subsection*{You'll want to know: der Ordner}
\textbf{der Ordner} — \textquotedblleft{}the folder; the ring binder\textquotedblright{}.

\begin{quote}
Alle Fotos sind in einem Ordner auf dem Computer.
\end{quote}

\begin{quote}
\emph{All the photos are in a folder on the computer.}
\end{quote}

Plural: \emph{die Ordner}. In an office, \emph{der Ordner} is a ring binder.

\begin{cousinweb}
\emph{Ordner} is formed with \emph{-er} on \emph{ordnen} \textquotedblleft{}to put in order\textquotedblright{}, inherited from Old High German \emph{ordinōn}, an early borrowing from Latin \emph{ōrdināre} (EtymDB), from \emph{ōrdo} \textquotedblleft{}row, order\textquotedblright{}.
\end{cousinweb}

\begin{grammarlens}[title={What you've built}]
One more word for telling, explaining and describing at B1.
\end{grammarlens}

\subsection*{Your turn}
\begin{itemize}
\raggedright
  \item \emph{Say it:} \emph{der Ordner}
  \item \emph{Say it:} \emph{der Ordner}, once more
  \item \emph{Say it:} a sentence of your own with \emph{der Ordner}
  \item \emph{Recall:} say the German for the, then the German for to save, then say \emph{der Ordner} again
\end{itemize}

\begin{tcolorbox}[breakable,colback=teal!4,colframe=teal!35!black,title={Before you move on}]
What does der Ordner mean? (\textquotedblleft{}The folder; the ring binder\textquotedblright{}.) Say it once more.
\end{tcolorbox}
\section[der Anhang]{der Anhang — the attachment; the appendix}

\label{lesson:GE-C393-anhang}



\begin{quote}
Before the new one: say the German for to save, then the German for the folder.
\end{quote}

\subsection*{You'll want to know: der Anhang}
\textbf{der Anhang} — \textquotedblleft{}the attachment; the appendix\textquotedblright{}.

\begin{quote}
Meinen Lebenslauf finden Sie im Anhang.
\end{quote}

\begin{quote}
\emph{You will find my CV attached.}
\end{quote}

Plural: \emph{die Anhänge}. In a book, the appendix.

\begin{cousinweb}
\emph{Anhang} goes with \emph{anhängen} \textquotedblleft{}to attach\textquotedblright{}, built on \emph{hängen} \textquotedblleft{}to hang\textquotedblright{}, Old High German \emph{hengen}, Proto-Germanic \emph{hangijaną} (EtymDB): what hangs on at the end.
\end{cousinweb}

\begin{grammarlens}[title={What you've built}]
One more word for telling, explaining and describing at B1.
\end{grammarlens}

\subsection*{Your turn}
\begin{itemize}
\raggedright
  \item \emph{Say it:} \emph{der Anhang}
  \item \emph{Say it:} \emph{der Anhang}, once more
  \item \emph{Say it:} a sentence of your own with \emph{der Anhang}
  \item \emph{Recall:} say the German for to save, then the German for the folder, then say \emph{der Anhang} again
\end{itemize}

\begin{tcolorbox}[breakable,colback=teal!4,colframe=teal!35!black,title={Before you move on}]
What does der Anhang mean? (\textquotedblleft{}The attachment; the appendix\textquotedblright{}.) Say it once more.
\end{tcolorbox}
\section[anhängen]{anhängen — to attach; to add on}

\label{lesson:GE-C393-anhaengen}



\begin{quote}
Before the new one: say the German for the folder, then the German for the attachment.
\end{quote}

\subsection*{You'll want to know: anhängen}
\textbf{anhängen} — \textquotedblleft{}to attach; to add on\textquotedblright{}.

\begin{quote}
Ich habe die Fotos an die E-Mail angehängt.
\end{quote}

\begin{quote}
\emph{I attached the photos to the email.}
\end{quote}

Separable: \emph{ich hänge an}, \emph{ich habe angehängt}.

\begin{cousinweb}
\emph{anhängen} is built from \emph{an} \textquotedblleft{}on\textquotedblright{} and \emph{hängen} \textquotedblleft{}to hang\textquotedblright{}, Old High German \emph{hengen}, Proto-Germanic \emph{hangijaną} \textquotedblleft{}to cause to hang\textquotedblright{} (EtymDB).
\end{cousinweb}

\begin{grammarlens}[title={What you've built}]
One more word for telling, explaining and describing at B1.
\end{grammarlens}

\subsection*{Your turn}
\begin{itemize}
\raggedright
  \item \emph{Say it:} \emph{anhängen}
  \item \emph{Say it:} \emph{anhängen}, once more
  \item \emph{Say it:} a sentence of your own with \emph{anhängen}
  \item \emph{Recall:} say the German for the folder, then the German for the attachment, then say \emph{anhängen} again
\end{itemize}

\begin{tcolorbox}[breakable,colback=teal!4,colframe=teal!35!black,title={Before you move on}]
What does anhängen mean? (\textquotedblleft{}To attach; to add on\textquotedblright{}.) Say it once more.
\end{tcolorbox}
\section[(dialogue)]{Words from chapters 386–393 — a review}

\label{lesson:GE-R393-review-386-393}



\begin{quote}
Say the German for to sue and the lawsuit.
\end{quote}

\subsection*{Your turn}
Say these aloud:

\begin{itemize}
\raggedright
  \item to sue, the lawsuit, the trial, to be liable, liability
  \item the bank balance, the credit, to overdraw, to debit, the amount
  \item the bank loan, the loan, the debts, to pay back, to borrow
  \item to afford, to be worth it, to finance, the expense, the income
  \item to waste, the waste, good value, the fortune, the payment
  \item the tax, the interest, to invest, to acquire, to insure
  \item the device, the rechargeable battery, to charge, the charger, the plug
  \item the, to save, the folder, the attachment, to attach
\end{itemize}

\begin{tcolorbox}[breakable,colback=teal!4,colframe=teal!35!black,title={Before you move on}]
To sue? (\textbf{klagen}.) The lawsuit? (\textbf{die Klage}.) The trial? (\textbf{der Prozess}.) To be liable? (\textbf{haften}.) Liability? (\textbf{die Haftung}.) The bank balance? (\textbf{der Kontostand}.) The credit? (\textbf{das Guthaben}.) To overdraw? (\textbf{überziehen}.) To debit? (\textbf{abbuchen}.) The amount? (\textbf{der Betrag}.) The bank loan? (\textbf{der Kredit}.) The loan? (\textbf{das Darlehen}.) The debts? (\textbf{die Schulden}.) To pay back? (\textbf{zurückzahlen}.) To borrow? (\textbf{borgen}.) To afford? (\textbf{leisten}.) To be worth it? (\textbf{lohnen}.) To finance? (\textbf{finanzieren}.) The expense? (\textbf{die Ausgabe}.) The income? (\textbf{die Einnahme}.) To waste? (\textbf{verschwenden}.) The waste? (\textbf{die Verschwendung}.) Good value? (\textbf{preiswert}.) The fortune? (\textbf{das Vermögen}.) The payment? (\textbf{die Zahlung}.) The tax? (\textbf{die Steuer}.) The interest? (\textbf{der Zins}.) To invest? (\textbf{investieren}.) To acquire? (\textbf{anschaffen}.) To insure? (\textbf{versichern}.) The device? (\textbf{das Gerät}.) The rechargeable battery? (\textbf{der Akku}.) To charge? (\textbf{aufladen}.) The charger? (\textbf{das Ladegerät}.) The plug? (\textbf{der Stecker}.) The? (\textbf{die Datei}.) To save? (\textbf{speichern}.) The folder? (\textbf{der Ordner}.) The attachment? (\textbf{der Anhang}.) To attach? (\textbf{anhängen}.)
\end{tcolorbox}
